%%%%%%%%%%%%%%%%%%%%% LATEX FILE %%%%%%%%%%%%%%%%%%%%%%%%%%%%%%%%%

\documentstyle[12pt]{article}
\textwidth 16cm
\textheight 24cm
\oddsidemargin 0cm
\evensidemargin 0cm
\topmargin -2cm
\headheight 0cm
\footheight 1cm

%theoremas
%\newtheorem{th}{Theorem}
%\newtheorem{prop}{Proposition}
\newtheorem{re}{Result}

% Make real R
%\font\Edth=msym10
%\def\Real{\hbox{\mid R}}

% Author
\Large
\begin{document}
\begin{center}
{\LARGE {\bf The characterization of third order}}\\ ~\\
{\LARGE {\bf ordinary differential equations}}\\ ~\\
{\LARGE {\bf admitting fiber-preserving point}}\\~\\ 
{\LARGE {\bf symmetries}} \\~\\
{\footnotesize Guy Grebot}\footnote{Supported by CNPq, Brazil.} \\
{\scriptsize Departamento de Matem\'{a}tica} \\
{\scriptsize Universidade de Bras\'{i}lia} \\
{\scriptsize Brasil.} \\
{\footnotesize Email:guy@mat.unb.br}
\end{center}

\begin{abstract}
Third order ordinary differential equations are characterized 
with respect to the dimension of the symmetry group of fiber preserving 
point transformations they admit.  For the symmetry groups of dimension 
greater or equal to 4, we obtain parametric families of equations 
admitting those groups. The equivalence classes of linear equations are 
also explicitly given. The results presented are 
based on the solution of the local equivalence problem for third order 
ordinary differential equations under fiber preserving point transformations.
\end{abstract}

\section{Introduction}

The local equivalence problem for third order ordinary differential 
equations (ODE's), consists in determining the conditions for 
a coordinate transformation, taking one third order ODE into 
another one, to exist. 
In \cite{T_Grebot94}, \'{E}lie Cartan's method of equivalence was  
used to solve this problem under the group $G$ of point transformations of 
the form $X=X(x)$, $Y=Y(x,y)$, where $y$ satisfies 
\begin{displaymath}
y'''=F(x,y,y',y''),
\end{displaymath}
the \lq $'$\rq denoting total differentiation with respect to $x$.

The method of equivalence directly gives the different equivalence classes 
which are characterized by a set of functions, invariant under a subgroup 
of a given prolongation of $G$. 
 
For self-equivalence, the group formed by the solutions 
of the problem is the symmetry group of the equation under study. Its 
dimension is related to the 
number of invariants determining the class of the considered equation 
(see \cite{P_Kamran89}).

The aim of this work is to give a characterization of third order 
ODE's admitting a symmetry group of fiber preserving point transformations.
This characterization is based on the possible values assumed by the 
invariants returned by the application of the equivalence method. 
The complexity of these expressions forbids us to give, as for second order 
ODE's \cite{P_HsuKam89}, 
the complete list of possible symmetry groups of third order ODE's. 
But when the invariants are assumed to be constant such a 
description is accessible, and we obtain a characterization of the 
third order ODE's admitting symmetry groups of dimension four to seven, as
well as a characterization of those groups.

The reason for not obtaining a characterization of the equations admitting 
lower dimensional symmetry groups is due to the high number of invariants. 
In fact, by the method of equivalence, a maximal functionally 
independent set must be selected among those invariants and their derivatives. 
This is always possible for 
particular cases with the help of the symbolic program {\tt clasode3} 
\cite{P_Grebot95}. But in the general case it is not possible: many 
possibilities appear as one must select at most three functionally 
independent invariants among, say, nine.

In the next section we briefly resume the solution of the 
equivalence problem for third order ODE's under fiber preserving point 
transformations. The linear equations are presented in 
Section \ref{s_lin}. In the other sections we give the characterization 
of the equations admitting symmetry groups of dimension 6 and 4.

In what follows, we have used the nomenclature of \cite{T_Grebot94}; the 
basic invariants, obtained by the application of the 
equivalence method, are named $I_{i}$ where $i \leq 10$. 


%%%%%%%%%%%%
\section{The solution of the equivalence problem}

We briefly present the solution, obtained in \cite{T_Grebot94}, to the problem 
of finding necessary and sufficient conditions for the existence of 
a transformation 
\begin{equation}
\label{point_trans}
\begin{array}{ccc}
X&=&\phi(x) \\
Y&=&\psi(x,y) 
\end{array} 
\end{equation}
taking one third order ODE 
\begin{equation}
\label{eq_F}
y'''=F(x,y,y',y''),
\end{equation}
into another one 
\begin{equation}
\label{eq_G}
Y'''=G(X,Y,Y',Y'').
\end{equation}
Such a problem is called an equivalence problem, and the method to solve it 
is E. Cartan's method of equivalence (see \cite{B_Gardn89}). 
We start with a coframe $\{ \omega \}$ 
\begin{eqnarray*}
\underline{\omega^{0}}&=& dx \\
\underline{\omega^{1}}&=& dy-p\,dx \\
\underline{\omega^{2}}&=& dp-q\,dx \\
\underline{\omega^{3}}&=& dq-F\,dx 
\end{eqnarray*}
defined on the second order jet bundle $J^{2}(R,R)$, with coordinates 
$(x,y,p,q)$, and a similar coframe $\underline{\Omega}$ determining the 
solutions of (\ref{eq_G}). 
The problem to be solved is to find the transformations 
$\Phi: J^{2}(R,R) \rightarrow J^{2}(R,R)$, which are given by the 
second prolongation of (\ref{point_trans}) such that 
\begin{equation}
\Phi^{\ast} \underline{\Omega}^{i}=\Gamma\,\underline{\omega}^{i},
\end{equation}
where $\Gamma$ are defined on $J^{2}(R,R)$ with values in the group 
$G$ of matrices of the form 
\begin{equation}
\label{Gamma}
\Gamma=\left(
\begin{array}{cccc}
A & 0 & 0 & 0 \\
0 & B & 0 & 0 \\
0 & C & B/A & 0 \\
0 & D & E & B/A^{2} 
\end{array} \right).
\end{equation}
The functions $A$, $B$, $C$, $D$ and $E$ are given by 
%%%%%%%%%%%%%%%%%%%%%%%%%%
\begin{eqnarray*} 
A&=&\phi_{,x} \\
B&=&\psi_{,y} \\
C&=&(\psi_{,xy}+p\,\psi_{,yy})/\phi_{,x} \\
D&=& \phi_{,x}^{-3}\,(q\,\phi_{,x}\,\psi_{,yy} 
+p\,(2\,\psi_{,xyy}\,\phi_{,x}-\psi_{,yy}\,\phi_{,xx}) \nonumber \\
& + & p^{2}\,\psi_{,yyy}\,\phi_{,x} 
+ \psi_{,xxy}\,\phi_{,x} - \phi_{,xx}\,\psi_{,xy}) \\
E&=& \phi_{,x}^{-3}\,(2\,\psi_{,xy}\,\phi_{,x}-\psi_{,y}\,\phi_{,xx}
+ 2\,p\,\psi_{,yy}\,\phi_{,x} ).
\end{eqnarray*}
%%%%%%%%%%%%%%%%%%%%%%%%%
The idea of the method is to reformulate the problem such that it 
becomes 
\begin{equation}
\tilde{\Phi}^{\ast} \Sigma^{i}=\sigma^{i}.
\end{equation}

The problem is completely solved in this form (see \cite{P_Kamran89}). 
The solution is given in terms of a set ${\cal F}$ of functions defined 
on $J^{2}(R,R)$. Among the elements of ${\cal F}$, a maximal set of 
functionally  independent functions must be selected. The dependence 
of the elements of ${\cal F}$ with respect to these independent functions 
 will characterize the equivalence class of a given equation. 

For self-equivalence, the solutions of the problem 
form the group of symmetries of the equation. The dimension $r$ of this group 
is related to the number $p$ of functionally independent invariants in 
${\cal F}$ 
and the dimension $n$ of the space on which the coframe $\{\sigma \}$ is 
defined. We have
\begin{equation}
r=n-p
\end{equation}
(see \cite{B_Gardn89} \cite{P_Kamran89} \cite{B_Sternberg64}).

For the solution of the equivalence problem of third order ODE's, 
the algorithm of the method splits depending on whether 
{\bf A}: $F_{,qq}=0$ or {\bf B}: $F_{,qq} \neq 0$. 

In the first case, $F_{,qq}=0$, we are led to consider two more cases 
{\bf A-a}: $TD_{0}=0$ and {\bf A-b}: $TD_{0} \neq 0$, where 
\begin{eqnarray}
\label{TD0}
TD_{0}&=&(18\,F_{,p}\,F_{,q} + 4\,F_{,q}^{3} 
- 18\,F_{,q}\,\frac{d}{dx}F_{,q} + 54\,F_{,y} 
- 27\,\frac{d}{dx}F_{,p} \nonumber \\
&+& 9\,\frac{d^{2}}{dx^{2}}F_{,q})/(54\,A^{3}).
\end{eqnarray}
If $TD_{0}=0$ then we obtain two more cases. These cases lead to the class 
of the trivial equation and to equivalence classes of equations admitting 
symmetry groups of dimension 3, 1 or 0. If $TD_{0} \neq 0$, then we 
must consider the cases
\begin{description}
\item[A-ba)] $TA_{1}$, $TA_{2}$, $TA_{3}$ not null,
\item[A-bb)] $TA_{1}=0$ and $TA_{2}$, $TA_{3}$ not null,
\item[A-bc)] $TA_{2}=0$ and $TA_{1}$, $TA_{3}$ not null,
\item[A-bd)] $TA_{3}=0$ and $TA_{1}$, $TA_{2}$ not null, 
\item[A-be)] $TA_{1}=TA_{2}=0$ and $TA_{3}$ not null,
\item[A-bf)] $TA_{1}=TA_{3}=0$ and $TA_{2}$ not null, 
\item[A-bg)] $TA_{2}=TA_{3}=0$ and $TA_{1}$ not null, 
\item[A-bh)] $TA_{1}$, $TA_{2}$, $TA_{3}$ all null,
\end{description}
where 
\begin{eqnarray}
\label{A_bTA1}
TA_{1}&=&( - 6\,A_{,p}\,F_{,q} + A_{,q}\,(9\,F_{,p} + F_{,q}^{2} 
- 3\,\frac{d}{dx}F_{,q}) + 18\,A_{,y})/(18\,B\,A) 
\nonumber \\
&+& ( - 18\,B\,A_{,p}\,A^{2}\,E + 9\,A_{,q}\,A^{4}\,E^{2})/(18\,B^{3}\,A) \\
\label{A_bTA2}
TA_{2}&=&(B\,A_{,p} - A_{,q}\,A^{2}\,E)/B^{2} \\
\label{A_bTA3}
TA_{3}&=&A_{,q}\,A/B,
\end{eqnarray}
and 
\begin{eqnarray}
\label{1A_bA3}
A^{3}&=&(18\,F_{,p}\,F_{,q} + 4\,F_{,q}^{3} - 18\,F_{,q}\,\frac{d}{dx}F_{,q}
+ 54\,F_{,y} - 27\,\frac{d}{dx}F_{,p} + 9\,\frac{d^{2}}{dx^{2}}F_{,q})/54. 
\nonumber \\ && 
\end{eqnarray}
All these cases, apart from {\bf A-be} and {\bf A-bf}, furnish a set of 
invariants 
that are used to determine classes of equivalence. The case {\bf A-bh} leads
 to the classes of the non trivial linear equations, and to other classes 
of equations admitting a 2-dimensional symmetry group. Apart from this case, 
the case {\bf A-ba} is the only one leading to possible 
classes of equations admitting a 4-dimensional group of symmetries. For that 
case, the possible 4-dimensional symmetry groups would appear by letting 
the invariants constant. In doing so, we observe that these constant do not 
satisfy the Jacobi identities, and define no Lie algebra.

In case {\bf A}, only the linear equations admit a maximal symmetry group 
of fiber preserving point transformations, of dimension 4,5 or 7.

In case {\bf B}: $F_{,qq} \neq 0$, we have to consider the following cases
\begin{description}
\item[B-a)] $TB_{3}=TB_{0}=TD_{0}=0$, 
\item[B-b)] $TB_{3}=TB_{0}=0$, $TD_{0}$ non null,
\item[B-c)] $TB_{3}=0$, $TB_{0}$ non null, $TB_{2}=-1$,
\item[B-d)] $TB_{3}=0$, $TB_{0}$ non null, $TB_{2} \neq -1$,
\item[B-e)] $TB_{3}$ non null,
\end{description}
where 
\begin{eqnarray}
\label{B_TB0}
TB_{0}&=&(F_{,qq}\,F_{,q} + 3\,\frac{d}{dx}\,F_{,qq})/(3\,A\,F_{,qq}) \\
%
\label{B_TB2}
TB_{2}&=&( - 9\,A\,E\,F_{,qqq} + 3\,F_{,pqq}\,F_{,qq})/F_{,qq}^{3} \\
%
\label{B_TB3}
TB_{3}&=&(3\,A\,F_{,qqq})/F_{,qq}^{2} \\
%
TD_{0}&=&(18\,F_{,p}\,F_{,q} 
+ 4\,F_{,q}^{3} - 18\,F_{,q}\,\frac{d}{dx}\,F_{,q} 
+ 54\,F_{,y} + 9\,\frac{d^{2}}{dx^{2}}\,F_{,q} \nonumber \\
&-& 27\,\frac{d}{dx}\,F_{,p})/(54\,A^{3}),
%
\end{eqnarray}
%
$A$ and $E$ being group parameters. Each of these cases will furnish 
a set of invariants which must be used to characterize the equivalence 
class of a given equation. Only case {\bf B-e} will allow equivalence classes 
of equations admitting a 4-dimensional symmetry group. The case 
{\bf B-a} leads to the equivalence class of the equations admitting a 
6-dimensional symmetry group. 

%%%%%%%%%%%%%%%%%%%%%%%%%%%%%%%%%%%%%%%%%%%%%%%%%%%%%%%%%%%%%%%%%%%%%%%%

\section{The Linear equations}
\label{s_lin}
In this section we give the charactrerization of the linear equations 
which admit fiber preserving symmetry groups of dimensions 4, 5 and 7.
 
Replacing $F$ in the expressions of the preceding section, by the expression 
of a canonical linear equation (see \cite{B_Gregus87})
\begin{equation}
\label{can_lin}
F=-2\,M(x)\,y' -(N(x)+M'(x))y
\end{equation}
we obtain the conditions determining the equivalence classes 
of the linear equations. There are two  possibilities for the 
so-called (see \cite{P_Chern40}) pseudo invariant $TD_{0}$ which assumes 
the form
\begin{equation}
TD_{0}=-\frac{N}{A^{3}}.
\end{equation}

>From this expression we directly see that the linear equations that 
can be equivalent to the trivial equation must satisfy $N=0$.  
This is exactly the condition found in the example 1 
of \cite{P_Boch1993}. 

The equivalence classes of the linear equations are then determined by the  
value of the so called Laguerre invariant $N$. 

The conditions characterizing the class of the trivial equation are given by
$TD_{0}=0$ and we obtain that a third order ordinary differential equation 
\begin{equation}
y'''=F(x,y,y',y'')
\end{equation}
can be transformed into the trivial equation 
\begin{equation}
y'''=0,
\end{equation}
by a transformation $X=\phi(x)$, $Y=\psi(x,y)$, if, and only if,
\begin{equation}
\label{1AAaaa}
\left\{
\begin{array}{rll}
F&=&p^{3}\,(1/3\,C_{0,y}-(C_{0}/3)^{2})+p^{2}\,(C_{0,x}-C_{0}/3\,C_{1}) \\
&+& p\,C_{2} + p\,q\,C_{0}+ q\,C_{1} + C_{3}\\
0&=&\frac{\partial}{\partial\,y}C_{3}-C_{3}\,C_{0}/3+TX/54 \\
TX&=&9\,\partial_{x}(C_{1,x}-C_{1}^{2}-3\,C_{2})
+C_{1}\,(18\,C_{2}+4\,C_{1}^{2}),
\end{array} \right.
\end{equation}
where $C_{i}=C_{i}(x,y)$.

The equations belonging to this class admit the maximal 7-dimensional group 
of symmetries. 

If $TD_{0} \neq 0$, then we obtain two other classes of linear equations which 
are characterized by the invariant
\begin{equation}
I_{1}=(A^{2}\,( - 3\,F_{6,xx} + F_{6,x}^{2} + 3\,F_{1}) + 6\,A\,A_{,xx} 
- 9\,A_{,x}^{2})/(6\,A^{4}),
\end{equation}
with $F_{i}=F_{i}(x,y)$ for $i=0,1,6$ and the expression for $A$ is given in 
(\ref{1A_bA3}).
The Maurer-Cartan equations for the symmetry group are given by 
\begin{equation}
\left\{ \begin{array}{lll}
d \omega^{0}&=&0 \\
d \omega^{1}&=&
\beta^{4}  \wedge  \omega^{1} 
+ \omega^{0}  \wedge  \omega^{2} \\
d \omega^{2}&=&
\beta^{4}  \wedge  \omega^{2} 
+ \omega^{0}  \wedge  \omega^{1}\,I_{1} + \omega^{0}  \wedge  \omega^{3} \\
d \omega^{3}&=&
\beta^{4}  \wedge  \omega^{3} 
+ \omega^{0}  \wedge  \omega^{1} + \omega^{0}  \wedge  \omega^{2}\,I_{1} \\
d \beta^{4}&=&0.
\end{array} \right.
\end{equation}
As $I_{1}$ depends only on $x$, there can be at most only one functionally 
independent invariant, meaning that we can either have a 5-dimensional or a 
4-dimensional symmetry group.

\begin{re}
There is a transformation $X=\phi(x)$, $Y=\psi(x,y)$, taking 
a third order ordinary differential equation 
\begin{equation}
y'''=F(x,y,y',y'')
\end{equation}
into the linear equation 
\begin{equation}
Y'''+2\;M\;Y'+N\;Y=0,
\end{equation}
where $M$ and $N$ are constant, if, and only if, 
\begin{eqnarray}
F&=&q\,(p\,F_{6,y} + F_{6,x}) + p^{3}\,(3\,F_{6,yy} - F_{6,y}^{2})/9 
\nonumber \\
&+& p^{2}\,(3\,F_{6,xy} - F_{6,x}\,F_{6,y})/3 + p\,F_{1} + F_{0},
\end{eqnarray}
where $F_{i}=F_{i}(x,y)$, and 
\begin{eqnarray}
F_{6,xxy}&=& (3\,F_{1,y} + 2\,F_{6,xy}\,F_{4})/3 \\
0&=&27\,F_{0,yy} - 9\,F_{0,y}\,F_{6,y} - 9\,F_{1,xy} + 3\,F_{1,y}\,F_{6,x} 
\nonumber \\
&-& 6\,F_{6,xx}\,F_{6,xy} + 2\,F_{6,xy}\,F_{6,x}^{2} + 9\,F_{6,xy}\,F_{1} 
- 9\,F_{6,yy}\,F_{0} \\
M/N^{2/3}&=&(A^{2}\,( - 3\,F_{6,xx} + F_{6,x}^{2} + 3\,F_{1}) 
+ 6\,A\,A_{,xx} - 9\,A_{,x}^{2})/(6\,A^{4}).
\end{eqnarray}
\end{re}
We see that the linear equations with constant coefficients are 
pa\-ra\-me\-tri\-zed by the constant $M/N^{2/3}$, which is the 
value of the invariant $I_{1}$ of case {\bf A-bha}. These equations admit 
a 5-dimensional symmetry group of fiber preserving point transformations. 

The characterization of the linear equations with non constant coefficients 
is given by the following result
\begin{re}
There is a transformation $X=\phi(x)$, $Y=\psi(x,y)$, taking 
a third order ordinary differential equation 
\begin{equation}
y'''=F(x,y,y',y'')
\end{equation}
into the linear equation 
\begin{equation}
Y'''+2\;M\;Y'+(M'+N)\;Y=0,
\end{equation}
where $M$ and $N$ are functions of the independent variable $X$ only, 
if, and only if, 
\begin{eqnarray}
F&=&q\,(p\,F_{6,y} + F_{6,x}) + p^{3}\,(3\,F_{6,yy} - F_{6,y}^{2})/9 
\nonumber \\
&+& p^{2}\,(3\,F_{6,xy} - F_{6,x}\,F_{6,y})/3 + p\,F_{1} + F_{0},
\end{eqnarray}
where $F_{i}=F_{i}(x,y)$, and 
\begin{eqnarray}
F_{6,xxy}&=& (3\,F_{1,y} + 2\,F_{6,xy}\,F_{4})/3 \\
0&=&27\,F_{0,yy} - 9\,F_{0,y}\,F_{6,y} - 9\,F_{1,xy} + 3\,F_{1,y}\,F_{6,x} 
\nonumber \\
&-& 6\,F_{6,xx}\,F_{6,xy} + 2\,F_{6,xy}\,F_{6,x}^{2} + 9\,F_{6,xy}\,F_{1} 
- 9\,F_{6,yy}\,F_{0},
\end{eqnarray}
and the equation
\begin{equation}
\begin{array}{lll}
& &-1/6\,N^{-8/3}\,(6\,M\,N^{2}+2\,N\,N_{,XX}-N_{,X}^{2}) \\
&=&(A^{2}\,( - 3\,F_{6,xx} + F_{6,x}^{2} + 3\,F_{1}) 
+ 6\,A\,A_{,xx} - 9\,A_{,x}^{2})/(6\,A^{4}).
\end{array}
\end{equation}
is satisfied.
\end{re}

Linear equations with non-constant coefficients will admit a 4-dimensional 
group of fiber-preserving point symmetries.


As a final remark, the equations given in the last two {\it results} 
furnish the sufficiency of the conditions for the existence of an 
equivalence, and can, in the case of the 
previous {\it results}, be used to determine the coordinate function $X$ 
of the transformation.

%%%%%%%%%%%%%%%%%%%%%%
\subsection{6-dimensional symmetry groups}
%SUBCASE B-a:
%%%%%%%%%%%%% Symmetry Analysis %%%%%%%%%%%%%%%%%
The only equations admitting a 6-dimensional symmetry group 
are of the form 
\begin{eqnarray*}
F&=&(12\,q^{2} + 36\,q\,F_{21,x} + 64\,p^{4}\,F_{01} 
+ 384\,p^{3}\,F_{21}\,F_{01} \nonumber \\
&+& 4\,p^{2}\,(216\,F_{21}^{2}\,F_{01} - F_{22}) 
+ 12\,p\,( - F_{21,xx} + 72\,F_{21}^{3}\,F_{01} - F_{21}\,F_{22}) 
\nonumber \\
&+& 9\,( - 2\,F_{21,xx}\,F_{21} + 3\,F_{21,x}^{2} + 36\,F_{21}^{4}\,F_{01} 
- F_{21}^{2}\,F_{22}))/(8\,p + 12\,F_{21}),
\end{eqnarray*}
where $F_{21}=F_{21}(x)$, $F_{22}=F_{22}(x)$ and $F_{01}=F_{01}(x,y)$, 
and the further condition
\begin{eqnarray*}
F_{01,x}=3/2\,F_{01,y}\,F_{21},
\end{eqnarray*}
%
has to be satisfied. All these equations are equivalent to 
\begin{equation}
\label{dim6_can}
y'''=\frac{3}{2}\frac{q^{2}}{p}.
\end{equation}
The symmetry group of fiber preserving point transformations for such 
equations is determined by the Maurer-Cartan equations 
%structure equations
%
\begin{equation}
\left\{
\begin{array}{lll}
d \omega^{0}&=&\alpha^{2} \wedge \omega^{0} \\
d \omega^{1}&=&\omega^{0} \wedge \omega^{2} + \omega^{1} \wedge \omega^{2} 
+ \alpha^{2} \wedge \omega^{1} \\
d \omega^{2}&=&\omega^{0} \wedge \omega^{3} - \omega^{1} \wedge \omega^{3} 
+ \mu^{2} \wedge \omega^{1} \\
d \omega^{3}&=& - \omega^{2} \wedge \omega^{3} + \mu^{2} \wedge \omega^{2} 
- \alpha^{2} \wedge \omega^{3} \\
d \alpha^{2}&=&2\,\omega^{0} \wedge \omega^{3} - \mu^{2} \wedge \omega^{0} \\
d \mu^{2}&=&2\,\omega^{2} \wedge \omega^{3} + \mu^{2} \wedge \alpha^{2} 
- 2\,\mu^{2} \wedge \omega^{2}.
\end{array} \right.
\end{equation}
%end structure equations
The transformations $X=\phi(x)$ and $Y=\psi(x,y)$, relating (\ref{dim6_can}) 
to another equation of the class are determined by
\begin{eqnarray*}
\psi_{,x}&=&\frac{3}{2}\,\psi_{,y}\,F_{21} \\
\psi_{,yyy}&=&\frac{3}{2}\frac{(\psi_{,yy})^{2}}{\psi_{,y}}-8\,F_{01}\,
\psi_{,y} \\
\phi'''&=&\frac{3}{2}\frac{\phi''^{2}}{\phi'}-\frac{\phi'}{2}F_{22}.
\end{eqnarray*}

%%%%%%%%%%%%% end of Symmetry Analysis %%%%%%%%%%%%%%%%%
%

%%%%%%%%%%%%%%% 4-dimensional %%%%%%%%%%%%%%%%%%%%%%%%
\subsection{4-dimensional symmetry groups}
The determination of the equations admitting a four dimensional 
symmetry group of fiber preserving point transformations is not as direct as  
it was for higher dimensions. In some cases, the result is direct, but  
in general, we have to assume all the invariants to be constant and look at  
the conditions imposed by this constraint. In doing so it is possible to 
show that the only equations admitting such groups are characterized as below.
%

As mentioned above, the only possible case for which a 4-dimensional symmetry 
group is possible is the case {\bf B-e}.
%
Let us assume that all the invariants given in equations 
(\ref{BeI1})-(\ref{BeI10}) in appendix \ref{ap_Be}, are constant, 
so that we can determine 
which equations have a 4-dimensional symmetry group of fiber preserving 
point transformations. $I_{4}$, $I_{3}$ lead to
\begin{eqnarray}
\frac{d}{dx}\,F_{,qqq}
&=&( - 3\,F_{,pqq} - 4\,F_{,qqq}\,F_{,q} - F_{,qq}^{2}\,I_{3}\,I_{6} 
+ F_{,qq}^{2}\,I_{6} - F_{,qq}^{2})/3 \\
%
F_{,qqqq}&=&F_{,qqq}^{2}\,(I_{3} + 2)/F_{,qq}.
\end{eqnarray}
$I_{6}$ and $I_{5}$ give
\begin{eqnarray}
\frac{d}{dx}\,F_{,qq}
&=&F_{,qq}\,( - F_{,qqq}\,F_{,q} + F_{,qq}^{2}\,I_{6})/(3\,F_{,qqq}) \\
%
\frac{d}{dx}\,F_{,q}&=&( - 9\,F_{,pqq}^{2} 
+ 6\,F_{,qqq}\,(F_{,pq}\,F_{,qq} + 3\,F_{,qqy} - F_{,pqq}\,F_{,q}) \nonumber \\
&+& F_{,qqq}^{2}\,(9\,F_{,p} + F_{,q}^{2}) - 6\,F_{,pqq}\,F_{,qq}^{2} 
+ 2\,F_{,qq}^{4}\,I_{5})/(3\,F_{,qqq}^{2}),
\end{eqnarray}
and from 
\mbox{$\partial_{q}\,\frac{d}{dx}\,F_{,q}-\frac{d}{dx}\,F_{,qq}
-F_{,pq}-F_{,q}\,F_{,qq}=0$} we get
\begin{eqnarray}
F_{,yqqq}&=&( - 6\,F_{,pqq}\,F_{,qq}^{2}\,I_{3} 
+ 6\,F_{,pq}\,F_{,qqq}\,F_{,qq}\,I_{3} + 18\,F_{,qqq}\,F_{,qqy}\,I_{3} 
+ 36\,F_{,qqq}\,F_{,qqy} \nonumber \\
&+& 4\,F_{,qq}^{4}\,I_{3}\,I_{5} + F_{,qq}^{4}\,I_{6})/(18\,F_{,qq}).
\end{eqnarray}
We can compute $\frac{d^{2}}{dx^{2}}\,F_{,q}$ and substitute it into the 
$I_{i}$. $I_{2}$ and $I_{7}$ respectively give
\begin{eqnarray}
F_{,pqqq}&=&(F_{,pqq}\,F_{,qqq}\,(I_{3} + 2))/F_{,qq} \\
%
F_{,ppqq}&=&(9\,F_{,pqq}^{2}\,I_{3} + 18\,F_{,pqq}^{2} 
+ 6\,F_{,pqq}\,F_{,qq}^{2} - 6\,F_{,pq}\,F_{,qqq}\,F_{,qq} 
- F_{,qq}^{4}\,I_{7} 
\nonumber \\&-& F_{,qq}^{4}\,I_{5})/(9\,F_{,qq}).
\end{eqnarray}
>From $I_{9}$ we obtain
\begin{eqnarray}
F_{,ppq}&=& - ((2\,I_{3}\,I_{6}^{2} + 8\,I_{3}\,I_{6}\,I_{7} + 4\,I_{3}\,I_{8} 
- I_{6}^{2} - 4\,I_{6}\,I_{5} \nonumber \\
&-& 8\,I_{6}\,I_{7} + 2\,I_{6} + 4\,I_{5} - 2\,I_{8})\,F_{,qq}^{4} \nonumber \\
&-& 18\,F_{,pqq}^{2} - 12\,F_{,pqq}\,F_{,qq}^{2} 
+ 12\,F_{,pq}\,F_{,qqq}\,F_{,qq})/(18\,F_{,qqq}).
\end{eqnarray}
$I_{8}$ furnishes
\begin{eqnarray}
F_{,ypqq}&=&( - 18\,F_{,pqq}^{2}\,F_{,qq}^{2}\,(2\,I_{3} - 3) 
+ 6\,F_{,pqq}\,(6\,(3\,(I_{3} + 2)\,F_{,qqy}+ F_{,qq}^{2}\,F_{,q})\,F_{,qqq} 
\nonumber \\&-& (4\,I_{3}\,I_{7} - I_{6} - 2\,I_{5} - 2\,I_{7} 
- 2)\,F_{,qq}^{4} + 6\,F_{,pq}\,F_{,qqq}\,F_{,qq}\,I_{3}) \nonumber \\
&+& F_{,qq}^{6}\,(4\,I_{3}\,I_{6}^{2}
+ 16\,I_{3}\,I_{6}\,I_{7} + 8\,I_{3}\,I_{8} - I_{6}^{2} - 12\,I_{6}\,I_{5} 
\nonumber \\&-& 16\,I_{6}\,I_{7}+ 4\,I_{6} - 4\,I_{7} - 2\,I_{8}) \nonumber \\
&-& 4\,F_{,qq}^{2}\,((I_{5} + I_{7})\,F_{,qq}^{2}\,F_{,q}
+ 9\,F_{,qqy})\,F_{,qqq} \nonumber \\
&+& 6\,F_{,qq}\,F_{,qqq}\,( - 2\,(3\,F_{,qqq}\,F_{,q} 
+ F_{,qq}^{2})\,F_{,pq} \nonumber \\
&-& 9\,F_{,pp}\,F_{,qqq} + 6\,F_{,qqq}\,F_{,yq}))/(108\,F_{,qq}\,F_{,qqq}).
\end{eqnarray}
The compatibility conditions between these expressions lead to 
\begin{eqnarray}
0&=&(I_{3}-1)\,(I_{5}+I_{7}) \\
0&=&(I_{3}-2)\,I_{9} \\
0&=&(I_{3}-1)\,(I_{6}+4\,I_{3}\,I_{5}),
\end{eqnarray}
and to 
\begin{eqnarray}
I_{1}&=& - (2\,I_{3}\,I_{7} - I_{6} - 2\,I_{5} - 2\,I_{7})/2 \\
I_{2}&=&0 \\
I_{4}&=&I_{3}\,I_{6} + 2\,I_{6} + 1 \\
I_{9}&=&(2\,I_{3}\,I_{6}^{2} + 8\,I_{3}\,I_{6}\,I_{7} + 4\,I_{3}\,I_{8} 
- I_{6}^{2} - 4\,I_{6}\,I_{5} \nonumber \\
&-& 8\,I_{6}\,I_{7} + 2\,I_{6} + 2\,I_{5} - 2\,I_{7} - 2\,I_{8})/2 \\
I_{10}&=& - (I_{3} - 1)\,I_{7}.
\end{eqnarray}

\noindent
For $I_{3}=1$, we obtain 
\begin{equation} \left\{
\begin{array}{lll}
I_{1}&=&(I_{6} + 2\,I_{5})/2 \\
I_{2}&=&0 \\
I_{4}&=&3\,I_{6} + 1 \\
I_{7}&=&1/2\,I_{6}^{2}+I_{8}-I_{6}\,(I_{5}-1)-I_{5}\,(I_{6}-1) \\
I_{9}&=&0 \\
I_{10}&=&0,
\end{array} \right.
\end{equation}
and $I_{5}$, $I_{6}$, $I_{8}$ are free constants.

\noindent
If $I_{3}=2$, then
\begin{equation} \left\{
\begin{array}{lll}
I_{1}&=& - 2\,I_{5} \\
I_{2}&=&0 \\
I_{4}&=& - 32\,I_{5} + 1 \\
I_{6}&=&-8\,I_{5} \\
I_{7}&=&-I_{5} \\
I_{9}&=&3\,(48\,I_{5}^{2} - 2\,I_{5} + I_{8}) \\
I_{10}&=&I_{5}, 
\end{array} \right.
\end{equation}
and $I_{5}$ and $I_{8}$ are free.

\noindent
For the other cases, i.e $I_{3} \neq 2$ and $I_{3} \neq 1$, we are left with 
\begin{equation} \left\{
\begin{array}{lll}
I_{7}&=&-I_{5} \\
I_{6}&=&-4\,I_{3}\,I_{5} \\
0&=&(2\,I_{3}-1)\,(8\,I_{3}\,I_{5}^{2}\,(I_{3}+1)-2\,I_{5}+I_{8}).
\end{array} \right.
\end{equation}
For $I_{3}=1/2$, $I_{5}$ and $I_{8}$ are free, and 
\begin{equation} \left\{
\begin{array}{lll}
I_{1}&=&( - I_{5})/2 \\
I_{2}&=&0 \\
I_{4}&=& - (5\,I_{5} - 1) \\
I_{6}&=& - 2\,I_{5} \\
I_{7}&=& - I_{5} \\
I_{9}&=&0 \\
I_{10}&=&( - I_{5})/2.
\end{array} \right.
\end{equation}
For $I_{3} \neq 1/2$, 
$I_{3}$ and $I_{5}$ are free and 
\begin{equation} \left\{
\begin{array}{lll}
I_{1}&=& - I_{3}\,I_{5} \\
I_{2}&=&0 \\
I_{4}&=& - (4\,I_{3}^{2}\,I_{5} + 8\,I_{3}\,I_{5} - 1) \\
I_{6}&=& - 4\,I_{3}\,I_{5} \\
I_{7}&=& - I_{5} \\
I_{8}&=&2\,I_{5}-8\,I_{3}\,I_{5}^{2}\,(I_{3}+1) \\
I_{9}&=&0 \\
I_{10}&=&(I_{3} - 1)\,I_{5}.
\end{array} \right.
\end{equation}
Replacing these values into the Maurer-Cartan equations 
\begin{equation} \left\{ \begin{array}{lll}
d \omega^{0}&=& \omega^{0} \wedge \omega^{1}\,I_{1} 
+ \omega^{0} \wedge \omega^{2}\,I_{2} 
+ \omega^{0} \wedge \omega^{3}\,I_{3} \\
d \omega^{1}&=&\omega^{0} \wedge \omega^{1}\,I_{4} 
+ \omega^{0} \wedge \omega^{2} 
+ \omega^{1} \wedge \omega^{2}\,I_{2} 
+ \omega^{1} \wedge \omega^{3}\,(I_{3} - 1) \\
d \omega^{2}&=&\omega^{0} \wedge \omega^{1}\,I_{5} 
+ \omega^{0} \wedge \omega^{2}\,I_{6} 
+ \omega^{0} \wedge \omega^{3} + \omega^{1} \wedge \omega^{2}\,I_{7} \\
&-& \omega^{1} \wedge \omega^{3}\,(I_{2} + 1) 
- \omega^{2} \wedge \omega^{3} \\
d \omega^{3}&=&\omega^{0} \wedge \omega^{1}\,I_{8} 
+ \omega^{0} \wedge \omega^{2}\,I_{5} 
+ \omega^{0} \wedge \omega^{3}\,(2\,I_{6} - I_{4}) \\
&+& \omega^{1} \wedge \omega^{2}\,I_{9} 
+ \omega^{1} \wedge \omega^{3}\,I_{10}
\end{array} \right.
\end{equation}
we obtain the characterization of the groups admitted by the different 
classes of equations, if the $I_{i}$ satisfy the Jacobi identities
\begin{eqnarray}
\label{jacobi}
\left\{
\begin{array}{rcl}
0&=& I_{3}\,I_{9} + I_{1}\,I_{2} + I_{2}\,I_{7} \\
0&=&  - I_{3}\,I_{1} - I_{3}\,I_{10} + I_{1} + I_{2}^{2}  + I_{2} \\
0&=& I_{2} \\
0&=& I_{3}\,I_{5} + I_{1} - I_{2}\,I_{4} + I_{2}\,I_{6} - I_{5} - I_{7} \\
0&=&  - 2\,I_{3}\,I_{4} + 2\,I_{3}\,I_{6} + 2\,I_{2} + I_{4} - 2\,I_{6} + 1 \\
0&=& I_{2} \\
0&=& I_{1}\,I_{6} - 3\,I_{2}\,I_{5} + I_{4}\,I_{7} - I_{5} + I_{8} - I_{9} \\
0&=&  - 2\,I_{3}\,I_{5} + I_{1} - I_{2}\,I_{6} - I_{6} + I_{7} - I_{10} \\
0&=&  - I_{3}\,I_{6} + I_{4} - 2\,I_{6} - 1 \\
0&=&  - I_{3}\,I_{7} - I_{2}^{2}  - I_{2} + I_{7} - I_{10} \\
0&=& I_{1}\,I_{5} - 2\,I_{2}\,I_{8} + 2\,I_{4}\,I_{9} 
- I_{5}\,I_{7} + I_{5}\,I_{10} - I_{6}\,I_{9} \\
0&=&  - 2\,I_{3}\,I_{8} - I_{1}\,I_{4} + 2\,I_{1}\,I_{6} + I_{2}\,I_{5} 
+ I_{4}\,I_{10} + I_{5} + I_{8} + I_{9} \\
0&=&  - I_{3}\,I_{5} - I_{2}\,I_{4} + 2\,I_{2}\,I_{6} + I_{5} + I_{10} \\
0&=&  - I_{3}\,I_{9} + I_{2}\,I_{10} + 2\,I_{9}.
\end{array} \right.
\end{eqnarray}

Replacing, for each of the cases above, the values obtained for the $I_{i}$, 
in (\ref{jacobi}), we obtain the results shown in table \ref{tab_1}. In the 
last column of this table, we have put the normal form of the commutators as 
furnished by the computer algebra program presented in \cite{P_Schobel92}. 
\begin{table}[h]
\label{tab_1}
\caption{4-dimensional Lie algebras for the case {\bf B-e}}
\begin{tabular}{|l|l|}
\hline \hline 
& \\
structure constants & normal form for the commutators \\
& and classification \\
\hline
& \\
& For $I_{5}=0$: \\
$I_{3}=1$, & $[X,Y]=W$, $[X,Z]=Y$, \\
$I_{1}=I_{7}=I_{5}$, &  $[W,X]=[W,Y]=[W,Z]=0$ \\
$I_{4}=1$, & $[Y,Z]=X$,  case 19 \\
$I_{2}=I_{6}=I_{8}=I_{9}=I_{10}=0$. & For $I_{5} \neq 0$: \\
& $[W,X]=Y$, $[W,Y]=X$, \\ 
&$[W,Z]=[X,Z]=[Y,Z]=0$ \\
& $[X,Y]=W$,  case 22 \\                   
\hline                 
& \\
%%%%%%%%%%%%%%%
$I_{3}=2$, & $[X,Y]=W$, $[X,Z]=4X$, \\
$I_{5}=\frac{1}{40}=\frac{-1}{2}I_{1}=\frac{1-I_{4}}{32}
=\frac{-I_{6}}{8}=-I_{7}=I_{10}$ & $[W,X]=[W,Y]=0$, \\
$I_{2}=0$, &  $[W,Z]=2W$, $[Y,Z]=-2Y$, \\
$I_{8}=\frac{1}{50}$, & case 16 with lambda=3 \\
$I_{9}=144\,I_{5}^{2} - 6\,I_{5} + 3\,I_{8}$. &  \\
& \\
\hline        
& \\           
%%%%%%%%%%%%%%%%%%%%%
$I_{3}=\frac{1}{2}$, & $[X,Y]=W$, $[X,Z]=4X$, \\
$I_{5}=\frac{ - 1}{2}=-2\,I_{1}=-2\,I_{10}
=\frac{1-I_{4}}{5}=\frac{-I_{6}}{2}=-I_{7}$ & $[W,X]=[W,Y]=0$, \\
$I_{2}=I_{9}=0$, & $[W,Z]=2W$ \\
$I_{8}=\frac{ - 5}{2}$. & $[Y,Z]=-2Y$, \\
& case 16 with lambda=3 \\
& \\
\hline                 
& \\ 
%%%%%%%%%%%%%%%%%%%%%
$I_{1}= - I_{3}\,I_{5}$,& \\ 
$I_{2}=0$, & \\ 
$I_{4}=-4\,I_{3}^{2}\,I_{5}-8\,I_{3}\,I_{5}+1$,& \\
$I_{6}=-4\,I_{3}\,I_{5}$, & \\
$I_{7}= - I_{5}$, & \\
$I_{8}=-8\,I_{3}^{2}\,I_{5}^{2}-8\,I_{3}\,I_{5}^{2}+2\,I_{5}$, & \\
$I_{9}=0$, & \\
$I_{10}=I_{3}\,I_{5} - I_{5}$, & No real Lie algebra. \\
$I_{3} \in \,\,\mid\!\!\!\!{\tt C}$ & \\
$I_{5}=\frac{I_{3} - 1}{2\,I_{3}^{2}\,(2\,I_{3} + 1)}$. & \\
& \\
\hline 
\end{tabular}
\end{table}
%%%%%%%%%%%%%%%%%%%%%%%%%%%%%%%%%%%%%%%%%%%%%%%%%%%%%%%%%%%
With the admissible values for $I_{i}$, we can go back to equations 
(\ref{BeI1})-(\ref{BeI10}) and get the conditions for the equations to 
admit a 4-dimensional symmetry group.

%%%%%%%%%%% end of Symmetry Analysis %%%%%%%%%%%%%%%%%%%%%%%%%
%
%%%%%%%%%%% end of 4-dimensional %%%%%%%%%%%%%%%%%%%%%%%%%%%%%%
%%%%%%%%%%%%%%%%%%%%%%%%%%%%%%%%%%%%%%%%%%%%%%%%%%%%%%%%%%%%%%%
%%%%%%%%%%%%%%%%%%%%%%%%%%%%%%%%%%%%%%%%%%%%%%%%%%%%%%%%%%%%%%%

\section{Conclusion}
The solution of the local equivalence problem for third order ODE's 
gives a complete characterization of the equations admitting symmetry 
groups of fiber preserving point transformations of dimension $\geq 4$. 

The different possible 4-dimensional groups are given as parametric 
families. For those groups, the use of a computer algebra program 
presented in \cite{P_Schobel92} allows us to classify their Lie 
algebras. In principal, it should be possible 
to obtain the symmetry generators from the information obtained. This 
avoids the solution of a system of equations. But in practice the author 
was not able to obtain these generators on a concrete example, as has been 
done in \cite{P_KarlMacC82}.

It is interesting to note that, unlike the case of second order ODE's, 
it is not possible to characterize the equations admitting symmetry groups 
of dimension less than 4, in the general case. Nonetheless, with the help 
of the computer algebra program {\tt clasode3}, several equations admitting 
such groups, such as some of the Painlev\'{e} equations of the third order, 
were clasified. 

%Bibliography
\begin{thebibliography}{99}

\bibitem{P_Boch1993} Bocharov A. V., Sokolov V. V., Svinolupov S. I. (1993) 
                     {\it On Some Equivalence Problems for Differential  
                     Equations}, preprint.
\bibitem{P_Chern40} Chern S.S. (1940) 
                   {\it The geometry of a differential equation 
                    $Y'''=F(X,Y,Y',Y'')$}, 
                   Tensor, N.S., Vol. 28, 173-176.
\bibitem{B_Gardn89} Gardner, R. (1989) 
                    {\it The Method of Equivalence and its Applications}, 
                    Siam, Philadelphia, Pensylvania.
\bibitem{T_Grebot94} Grebot G. (1994) 
                     {\it On Killing tensors in vacuum Space-times and 
                     The classification of third order  
                      ordinary differential equations},
                     PhD Thesis, Queen Mary College, London.
\bibitem{P_Grebot95} Grebot G. (1995) 
                    {\it The classification of third order ordinary 
                    differential equations},
                    Proceedings of the CATHODE workshop, Nijmegen, 
                    9-12 January 1995, p. 23.
\bibitem{B_Gregus87} Gregus, M. (1987) 
                     {\it Third Order Linear Differential Equations}, 
                     Mathematics and its Applications, 
                     D. Reidel Publishing Company, Dordrecht, Holland.
\bibitem{P_HsuKam89} Hsu L., Kamran N. (1989) 
                     {\it Classification of Second-Order Differential 
                     Equations admitting Lie Groups of Fiber-preserving 
                     point Symmetries}, 
                     Proc. Lond. Math. Soc. (3), {\bf 58}, No 2, 387-416.
\bibitem{P_Kamran89} Kamran, N. (1989) 
                     {\it Contributions to the Study of the Equivalence 
                     problem of Elie Cartan and its Applications 
                     to partial and Ordinary Differential Equations}, 
                     Memoires Cl. Sci. Acad. Roy. Belg., {\bf 45}, Fac.7.
\bibitem{P_KamShd86} Kamran N., Shadwick W.F. (1986) 
                     {\it The solution of the Cartan Equivalence Problem for 
                      $\frac{d^{2}\,y}{dx^{2}}=F(x,y,\frac{d\,y}{dx})$ 
                      under the Pseudo-group $\overline{x}=\phi(x)$, 
                      $\overline{y}=\psi(x,y)$}, 
                      Lecture Notes in Physics {\bf 246}, 320-336.
\bibitem{P_KarlMacC82} Karlhede A. and MacCallum M.A.H. (1982) 
                       {\it On Determining the Isometry Group of a Riemannian 
                       Space}, 
                       Gen. Rel. Grav., Vol 14, No 7, 673.
\bibitem{P_Schobel92} F. Sch\"{o}bel (1992)
                   {\it The Symbolic Classification of Real Four Dimensional
                    Lie Algebras},
                    Universit\"{a}t Leipzig Preprint.
\bibitem{B_Sternberg64} Sternberg, S. (1964) 
                     {\it Lectures on differential geometry}, 
                    Prentice Hall, Englewood Cliffs, NJ.
\end{thebibliography}

\begin{appendix}
\section{Invariants for the case B-e}
\label{ap_Be}
The invariants obtained for the case {\bf B-e} are given below.
\begin{eqnarray}
\label{BeI1}
I_{1}&=&-(18\,F_{,pqqq}\,F_{,pqq}\,F_{,qqq}\,F_{,qq} + 6\,F_{,pqqq}
\,F_{,qqq}^{2}\,F_{,qq}\,F_{,q} \nonumber \\
&-& 9\,F_{,pqq}^{2}\,F_{,qqqq}\,F_{,qq} - 18\,F_{,pqq}^{2}\,F_{,qqq}^{2} 
- 12\,F_{,pqq}\,F_{,qqq}^{3}\,F_{,q} \nonumber \\
&-& 9\,F_{,p}\,F_{,qqqq}\,F_{,qqq}^{2}\,F_{,qq} 
+18\,F_{,p}\,F_{,qqq}^{4} - F_{,qqqq}\,F_{,qqq}^{2}\,F_{,qq}\,F_{,q}^{2} 
\nonumber \\&
+& 3\,F_{,qqqq}\,F_{,qqq}^{2}\,F_{,qq}\,\frac{d}{dx}\,F_{,q}
- 18\,F_{,yqqq}\,F_{,qqq}^{2}\,F_{,qq} + 2\,F_{,qqq}^{4}\,F_{,q}^{2} 
\nonumber \\&-& 6\,F_{,qqq}^{4}\,\frac{d}{dx}\,F_{,q} 
+ 36\,F_{,qqq}^{3}\,F_{,yqq})/(2\,F_{,qqq}^{2}\,F_{,qq}^{4}) \\
%
\label{BeI2}
I_{2}&=&3\,( F_{,pqqq}\,F_{,qqq} - F_{,pqq}\,F_{,qqqq})
/(F_{,qqq}^{2}\,F_{,qq}) \\
%
\label{BeI3}
I_{3}&=&( F_{,qqqq}\,F_{,qq} - 2\,F_{,qqq}^{2})/F_{,qqq}^{2} \\
%
\label{BeI4}
I_{4}&=&( - 3\,F_{,pqq}\,F_{,qq} - F_{,qqq}\,F_{,qq}\,F_{,q} 
+ 9\,F_{,qqq}\,\frac{d}{dx}\,F_{,qq} \nonumber \\
&-& 3\,F_{,qq}\,\frac{d}{dx}\,F_{,qqq})/F_{,qq}^{3} \\
%
\label{BeI5}
I_{5}&=&( - 9\,F_{,pqq}^{2} - 12\,F_{,pqq}\,F_{,qqq}\,F_{,q} 
- 18\,F_{,pqq}\,\frac{d}{dx}\,F_{,qqq} \nonumber \\
&-& 9\,F_{,p}\,F_{,qqq}^{2} - F_{,qqq}^{2}\,F_{,q}^{2} 
+ 3\,F_{,qqq}^{2}\,\frac{d}{dx}\,F_{,q} \nonumber \\
&-& 18\,F_{,qqq}\,F_{,yqq} 
+ 18\,F_{,qqq}\,\partial_{p}\,\frac{d}{dx}\,F_{,qq})/(2\,F_{,qq}^{4}) \\
%
\label{BeI6}
I_{6}&=&F_{,qqq}\,(F_{,qq}\,F_{,q} + 3\,\frac{d}{dx}\,F_{,qq})/F_{,qq}^{3} \\
%
\label{BeI7}
I_{7}&=&-(18\,F_{,ppqq}\,F_{,qqq}^{2}\,F_{,qq} 
- 36\,F_{,pqqq}\,F_{,pqq}\,F_{,qqq}\,F_{,qq} \nonumber \\
&+& 18\,F_{,pqq}^{2}\,F_{,qqqq}\,F_{,qq} 
+ 9\,F_{,pqq}^{2}\,F_{,qqq}^{2} + 6\,F_{,pqq}\,F_{,qqq}^{3}\,F_{,q} 
\nonumber \\&-& 6\,F_{,pqq}\,F_{,qqq}^{2}\,F_{,qq}^{2} 
+ 6\,F_{,pq}\,F_{,qqq}^{3}\,F_{,qq} - 9\,F_{,p}\,F_{,qqq}^{4} \nonumber \\
&-& F_{,qqq}^{4}\,F_{,q}^{2} 
+3\,F_{,qqq}^{4}\,\frac{d}{dx}\,F_{,q} - 18\,F_{,qqq}^{3}\,F_{,yqq})
/(2\,F_{,qqq}^{2}\,F_{,qq}^{4}) \\
%
\label{BeI8}
I_{8}&=&F_{,qqq}^{3}\,(18\,F_{,p}\,F_{,q} + 4\,F_{,q}^{3} 
- 18\,F_{,q}\,\frac{d}{dx}\,F_{,q} 
\nonumber \\&+& 54\,F_{,y} + 9\,\frac{d^{2}}{dx^{2}}\,F_{,q} 
- 27\,\frac{d}{dx}\,F_{,p})/(2\,F_{,qq}^{6}) \\
%
\label{BeI9}
I_{9}&=&3\,( - 18\,F_{,ppqq}\,F_{,pqq}\,F_{,qqq}^{2} 
- 6\,F_{,ppqq}\,F_{,qqq}^{3}\,F_{,q} + 9\,F_{,pp}\,F_{,qqq}^{4} 
\nonumber \\&+& 27\,F_{,pqqq}\,F_{,pqq}^{2}\,F_{,qqq} 
+ 6\,F_{,pqqq}\,F_{,pqq}\,F_{,qqq}^{2}\,F_{,q} 
\nonumber \\&+& 9\,F_{,pqqq}\,F_{,p}\,F_{,qqq}^{3} 
+F_{,pqqq}\,F_{,qqq}^{3}\,F_{,q}^{2} 
- 3\,F_{,pqqq}\,F_{,qqq}^{3}\,\frac{d}{dx}\,F_{,q} 
\nonumber \\&+& 18\,F_{,ypqq}\,F_{,qqq}^{3} 
- 9\,F_{,pqq}^{3}\,F_{,qqqq} + 6\,F_{,pqq}^{2}\,F_{,qqq}^{2}\,F_{,qq} 
\nonumber \\&-&12\,F_{,pqq}\,F_{,pq}\,F_{,qqq}^{3} 
- 9\,F_{,pqq}\,F_{,p}\,F_{,qqqq}\,F_{,qqq}^{2} \nonumber \\
&-& F_{,pqq}\,F_{,qqqq}\,F_{,qqq}^{2}\,F_{,q}^{2} 
+ 3\,F_{,pqq}\,F_{,qqqq}\,F_{,qqq}^{2}\,\frac{d}{dx}\,F_{,q} \nonumber \\
&-& 18\,F_{,pqq}\,F_{,yqqq}\,F_{,qqq}^{2} 
+ F_{,pqq}\,F_{,qqq}^{3}\,F_{,qq}\,F_{,q} \nonumber \\
&+& 3\,F_{,pqq}\,F_{,qqq}^{3}\,\frac{d}{dx}\,F_{,qq} 
+ 2\,F_{,pq}\,F_{,qqq}^{4}\,F_{,q} 
- 3\,F_{,qqq}^{4}\,\partial_{p}\,\frac{d}{dx}\,F_{,q}) \nonumber \\
&/&(2\,F_{,qqq}^{2}\,F_{,qq}^{5}) \\
%
\label{BeI10}
I_{10}&=&( - 9\,F_{,pqq}^{2} - 6\,F_{,pqq}\,F_{,qqq}\,F_{,q}
- 6\,F_{,pqq}\,F_{,qq}^{2} \nonumber \\
&+& 6\,F_{,pq}\,F_{,qqq}\,F_{,qq} + 9\,F_{,p}\,F_{,qqq}^{2} 
+ F_{,qqq}^{2}\,F_{,q}^{2} \nonumber \\
&-& 3\,F_{,qqq}^{2}\,\frac{d}{dx}\,F_{,q} + 18\,F_{,qqq}\,F_{,yqq} 
- F_{,qqq}\,F_{,qq}^{2}\,F_{,q} \nonumber \\
&-&3\,F_{,qqq}\,F_{,qq}\,\frac{d}{dx}\,F_{,qq})/(2\,F_{,qq}^{4})+I_{1}.
\end{eqnarray}
\end{appendix}
\end{document}
